\section{Программная реализация прототипа системы обнаружения}
\label{ch:impl}

\noindent
Содержание главы соответствует третьему этапу постановки задачи:
разработка программного прототипа системы обнаружения аномалий
(проектирование архитектуры, реализация модулей предобработки и
формирования окон, ядра детекции, минимального интерфейса),
обучение и оптимизация моделей с подготовкой реальных данных,
разработка конечно-автоматного агента и библиотеки атакующих
сценариев.

\subsection{Стек и общая структура программного проекта}

Прототип реализован на платформе Windows 11 в среде Visual Studio
Code в виде Python-пакета \texttt{diploma\_nids}. Технологический
стек выбран по результатам сравнительного анализа в открытых
публикациях.

\begin{itemize}
    \item Python 3.10+ (тестировался на 3.13.11) --- основной язык.
    \item PyTorch 2.x (CPU-сборка) --- глубокое обучение; выбран по
    зрелости детерминированных CPU-ядер и доминирующему положению в
    публикациях NIDS 2024--2025.
    \item scikit-learn 1.x --- классические baseline-модели.
    \item XGBoost --- gradient boosting baseline~\cite{chen2016xgboost}.
    \item Pydantic v2 --- валидация конфигов и схем данных.
    \item PyYAML --- параметризация всех стадий pipeline.
    \item FastAPI + uvicorn --- REST API инференса.
    \item Streamlit --- визуальная панель мониторинга.
    \item pytest --- юнит-тестирование (11 модулей).
    \item Make --- оркестрация сборки и воспроизводимости.
\end{itemize}

Общая структура проекта приведена в таблице~\ref{tab:project-layout}.
Общий объем исходного кода --- порядка 4000 строк Python.

{
\setlength{\tabcolsep}{4pt}
\small
\begin{longtable}{|p{2.5cm}|p{4.5cm}|p{6.8cm}|}
\caption{Структура программного проекта \texttt{diploma\_nids}}
\label{tab:project-layout}\\
\hline
Каталог & Назначение & Ключевые модули \\
\hline
\endfirsthead
\multicolumn{3}{@{}l@{}}{Продолжение таблицы~\ref{tab:project-layout}}\\[2pt]
\hline
Каталог & Назначение & Ключевые модули \\
\hline
\endhead
\hline
\endfoot
configs/      & 21 YAML-файл для всех стадий &
data, preprocess, models, train, eval, attacker, pipeline \\
\hline
data/         & Схемы, загрузка, препроцессинг, окна, splits &
schema.py, load.py, preprocess.py, windowing.py, splits.py \\
\hline
models/       & Модели-кандидаты, единый registry &
base.py, cnn\_lstm.py, rnn\_family.py, cnn1d.py, tcn.py,
transformer.py, autoencoder.py, classical.py \\
\hline
training/     & Обучение DL-моделей &
trainer.py, losses.py (focal, BCE) \\
\hline
eval/         & Метрики, калибровка, drift &
metrics.py, thresholding.py, drift.py, error\_analysis.py \\
\hline
attacker/     & Конечно-автоматный агент &
templates.py, agent.py, drift\_injector.py, runtime.py \\
\hline
inference/    & Online-инференс и алертинг &
stream.py, alerting.py, service.py (FastAPI) \\
\hline
ui/           & Streamlit-дэшборд &
streamlit\_app.py \\
\hline
utils/        & Воспроизводимость и I/O &
seed.py, io.py, logging.py \\
\hline
scripts/      & Точки входа CLI & 15 нумерованных скриптов:
загрузка UNSW, препроцессинг, обучение, оценка, прогон агента,
drift-тест, demo-стенд, эксперименты E1--E6 \\
\hline
tests/        & Юнит-тесты (11 модулей) & 86 проходящих тестов \\
\hline
\end{longtable}
}

\subsection{Подсистема обработки данных}

Схема UNSW-NB15 формализована Pydantic-моделями (модуль
\texttt{data/schema.py}). Загрузчик \texttt{data/load.py} валидирует
наличие признаков и меток, проверяет существование официального
train/test-разделения и возвращает \texttt{DataFrame}. Аналогично
реализована загрузка CICIDS2017 для cross-evaluation.

Класс \texttt{Preprocessor} реализует sklearn-совместимый интерфейс
\texttt{fit/transform/fit\_transform/save/load} с полной JSON-сериализацией
фитированного состояния. На стадии \texttt{fit} формируются:
квантильные границы клиппинга (0,001 и 0,999);
статистики robust-scaler (медиана, IQR);
one-hot категории для признаков с малой кардинальностью и
frequency-encoding для большой кардинальности;
фиксированный порядок выходных столбцов \texttt{output\_columns},
обеспечивающий идентичность представления в train/val/test/inference.

Модуль \texttt{data/windowing.py} реализует векторизованный
\texttt{build\_windows} с тремя стратегиями агрегации меток
(\texttt{last}, \texttt{majority}, \texttt{any}). Применяется
\texttt{np.lib.stride\_tricks.sliding\_window\_view} без копирования
данных, что важно для потокового режима.

\subsection{Реализованные модели}

В таблице~\ref{tab:models-impl-3} приведено 15 реализованных
моделей. Регистрация выполняется при импорте подпакета
\texttt{models} через декоратор \texttt{@register("name")}.

\begin{table}[H]
\caption{Реализованные модели в \texttt{diploma\_nids/models/}}
\label{tab:models-impl-3}
\centering
\small
\begin{tabularx}{\linewidth}{|L{2.5cm}|L{3.4cm}|c|Z|}
\hline
Имя & Тип & Окна & Файл \\
\hline
mlp                  & DL supervised             & нет & mlp.py \\
\hline
cnn1d                & DL supervised             & да  & cnn1d.py \\
\hline
cnn\_lstm            & DL supervised (proposed)  & да  & cnn\_lstm.py \\
\hline
lstm                 & DL supervised             & да  & rnn\_family.py \\
\hline
gru                  & DL supervised             & да  & rnn\_family.py \\
\hline
bilstm               & DL supervised + attn      & да  & rnn\_family.py \\
\hline
tcn                  & DL supervised             & да  & tcn.py \\
\hline
transformer          & DL supervised             & да  & transformer.py \\
\hline
autoencoder          & DL semi-supervised        & нет & autoencoder.py \\
\hline
vae                  & DL semi-supervised        & нет & autoencoder.py \\
\hline
\makecell[l]{logistic\_\\ regression} & classical                 & нет & classical.py \\
\hline
\makecell[l]{random\_\\ forest}       & classical                 & нет & classical.py \\
\hline
xgboost              & classical                 & нет & classical.py \\
\hline
\makecell[l]{isolation\_\\ forest}    & classical (unsup.)        & нет & classical.py \\
\hline
ocsvm                & classical (one-class)     & нет & classical.py \\
\hline
\end{tabularx}
\end{table}

Все DL-модели реализованы на PyTorch и единообразно возвращают
объект \texttt{ModelOutput} с полями \texttt{logits},
\texttt{probs} и \texttt{extras}. Архитектура CNN-LSTM:
\begin{equation*}
\begin{aligned}
\text{Input}(B, W, F)\ &\to\ [\text{Conv1D}\!\to\!\text{BN}\!\to\!\text{ReLU}\!\to\!\text{MaxPool}\!\to\!\text{Dropout}]\!\times\!2\\
&\to\ \text{BiLSTM}\ \to\ \text{Attention\,Pool}\ \to\ \text{LayerNorm}\ \to\ \text{MLP-head}\ \\ \to\ \text{Sigmoid}.
\end{aligned}
\end{equation*}
Число параметров --- порядка 200 тыс., размер модели на диске ---
около 0,8 МБ.

\subsection{Подсистема обучения и оценки}

Loss-функции. В \texttt{training/losses.py} реализованы binary
cross-entropy с логитами и фокальная функция
потерь~\cite{lin2017focal} с параметрами $\alpha$, $\gamma$,
задаваемыми в \texttt{configs/train/default.yaml}.

Тренер. \texttt{training/trainer.py} реализует AdamW + cosine
LR-scheduler, gradient clipping, early stopping по
\texttt{val\_pr\_auc}, сохранение лучшего чекпойнта и журнала
истории обучения.

Метрики. \texttt{eval/metrics.py} реализует Precision, Recall, F1,
FPR, MCC, ROC-AUC, PR-AUC, ECE, FP-per-time, TTD, поиск порога по
целевому FPR и по F1-max, bootstrap-CI.

Калибровка. \texttt{eval/thresholding.py} реализует температурное
масштабирование (формула~\ref{eq:temp-scaling}) через L-BFGS на
валидационных логитах.

Drift-метрики. \texttt{eval/drift.py} реализует PSI, KL и MMD с
автобандвидтом по медианной эвристике; функция
\texttt{drift\_report} возвращает агрегат с явным индикатором
\texttt{drift\_alarm} при превышении PSI-порога 0{,}25.

\subsection{Конечно-автоматный агент}

\subsubsection{Шаблоны атак}

Реализованы семь параметризованных шаблонов
(таблица~\ref{tab:attack-templates-3}), каждый формирует
последовательность объектов \texttt{FlowRecord} с явной привязкой к
классу \texttt{attack\_cat}.

\begin{table}[H]
\caption{Реализованная библиотека шаблонов атак}
\label{tab:attack-templates-3}
\centering
\small
\begin{tabularx}{\linewidth}{|L{2.6cm}|Z|L{2.6cm}|}
\hline
Шаблон & Параметризация & Класс UNSW \\
\hline
\makecell[l]{ddos\_\\ synflood} &
интенсивность $\lambda_0$, длительность $\Delta$, доля SYN-флагов &
DoS \\
\hline
\makecell[l]{port\_\\ scan} &
скорость подключений $\rho$, диапазон портов, режим &
Reconnaissance \\
\hline
\makecell[l]{brute\_\\ force} &
интервал между попытками, целевой порт, доля успехов &
Reconnaissance / Exploits \\
\hline
\makecell[l]{exploit\_\\ payload} &
размер payload, сочетания TCP-флагов, retransmissions &
Exploits \\
\hline
\makecell[l]{data\_\\ exfil} &
длительность сессии, асимметрия $b_{out}/b_{in}$ &
Backdoors \\
\hline
\makecell[l]{internal\_\\ recon} &
фан-аут потоков, периодичность, размер сессии &
Reconnaissance \\
\hline
\makecell[l]{worm\_\\ lateral} &
число направленных потоков, периодичность, размер сессии &
Worms \\
\hline
\end{tabularx}
\end{table}

Реализация дуальная: эмпирический режим
(\texttt{build\_templates\_from\_} \\ \texttt{dataframe}) фитится на выборке
UNSW-NB15-train и эмитит потоки по empirical-CDF числовых признаков
с jitter; параметрический режим
(\texttt{default\_template\_registry}) используется как fallback.

\subsubsection{Конечно-автоматный планировщик}

\texttt{attacker/agent.py} реализует \texttt{FSMAgent} с поддержкой
11 состояний (NORMAL, 7 ATTACK и 3 DRIFT), вероятностями переходов
из \texttt{configs/attacker/policy.yaml} и минимальными
длительностями состояний. Каждый тик автомата возвращает структуру
\texttt{AgentTick} с метаданными ground truth
(\texttt{fsm\_state}, \texttt{attack\_cat}, \texttt{drift\_type}),
необходимыми для расчета per-state recall в Эксперименте~E4.

\subsubsection{Инжектор дрейфа и runtime}

\texttt{attacker/drift\_injector.py} реализует три типа дрейфа:
covariate (сдвиг масштаба и шума), prior (изменение долей классов
через политику FSM), concept (маскирование сигнатурных признаков
атак).

\texttt{attacker/runtime.py} объединяет агент и drift-инжектор и
предоставляет два режима работы: offline (метод
\texttt{collect\_dataframe}) для эксперимента E4 и online
(\texttt{stream\_with\_callback}) для realtime-демо.

\subsection{Подсистема online-инференса и интерфейс}

\subsubsection{FastAPI-сервис}

Модуль \texttt{inference/service.py} реализует REST API:

\begin{itemize}
    \item \texttt{GET /health} --- проверка живости;
    \item \texttt{GET /info} --- сведения о загруженной модели;
    \item \texttt{POST /score} --- скоринг батча flow-записей;
    \item \texttt{POST /score-window} --- прямой скоринг готового
    окна $(W, F)$;
    \item \texttt{GET /alerts/recent} --- последние $n$ алертов.
\end{itemize}

Загрузка артефактов происходит через переменные окружения
\texttt{DIPLOMA\_MODEL\_YAML}, \texttt{DIPLOMA\_CHECKPOINT},
\texttt{DIPLOMA\_PREPROCESSOR}, \texttt{DIPLOMA\_THRESHOLD}, что
обеспечивает декларативность конфигурации.

\subsubsection{Формирование алертов}

\texttt{AlertFormer} реализует пятиступенчатую шкалу severity
(info, low, medium, high, critical), привязанную к нормированному
превышению скоринга над порогом, и временное окно дедупликации
($\Delta t = 5$\,с) для подавления шумовых повторов.

\subsubsection{Streamlit-дэшборд}

\texttt{ui/streamlit\_app.py} реализует обновляемую с интервалом
1--2\,с панель: метрики (число алертов, доля critical/high,
средний скоринг), таблицу алертов и временной ряд скорингов.
Сообщается с FastAPI-сервисом через \texttt{DIPLOMA\_API\_URL}.

\subsection{Воспроизводимость и журналирование}

\begin{itemize}
    \item Seed фиксируется через \texttt{utils/seed.py}
    (\texttt{numpy}, \texttt{random}, \texttt{torch},
    \texttt{torch.cuda}, \texttt{PYTHONHASHSEED}).
    \item Артефакты обучения сохраняются в JSON-формате в
    \texttt{experiments/runs/}; чекпойнты --- в \texttt{models/}.
    \item Конфигурации экспериментов фиксируются YAML-файлами;
    \texttt{make reproduce} запускает полный цикл
    data $\to$ preprocess $\to$ train $\to$ evaluate одной
    командой.
    \item Препроцессор сохраняется JSON-ом и применяется идентично
    в обучении и инференсе, исключая test-time leakage.
\end{itemize}

\subsection{Покрытие тестами}

Реализованы 11 модулей юнит-тестов (pytest), фиксирующих
критические инварианты программного прототипа; на текущей среде
проходят все 86 тестов. Состав модулей и проверяемые свойства
приведены в таблице~\ref{tab:tests}.

\begin{table}[H]
\caption{Состав юнит-тестов программного прототипа}
\label{tab:tests}
\centering
\small
\begin{tabularx}{\linewidth}{|L{3.6cm}|Z|}
\hline
Файл теста & Проверяемые свойства \\
\hline
test\_utils.py        & seed-воспроизводимость; YAML/JSON roundtrip;
создание родительских директорий; numpy в JSON \\
\hline
test\_configs.py      & все YAML-конфиги загружаются; FSM содержит
11 состояний; вероятности переходов суммируются в 1{,}0 \\
\hline
test\_preprocess.py   & shape выходов; отсутствие NaN/Inf;
неизменность статистик при повторном fit; save/load roundtrip \\
\hline
test\_windowing.py    & корректные shapes; три стратегии агрегации
меток; защита от коротких последовательностей; валидация входов \\
\hline
test\_splits.py       & сохранение временного порядка;
воспроизводимость random; непересечение split-ов \\
\hline
test\_models\_smoke.py& все 15 моделей загружаются из YAML и
выполняют forward; save/load CNN-LSTM; proposed = true \\
\hline
test\_losses.py       & FocalLoss положительный; correct < wrong;
build\_loss dispatch \\
\hline
test\_metrics.py      & Precision, Recall, F1, MCC; threshold для
target-FPR; ECE снижается калибровкой; bootstrap CI \\
\hline
test\_drift.py        & PSI/KL/MMD на одинаковых и сдвинутых
распределениях; drift\_report срабатывает на сдвиге \\
\hline
test\_attacker.py     & генерация всех 7 шаблонов; FSM эмиссия
тиков; drift-инжектор сдвигает признаки; runtime воспроизводим \\
\hline
test\_alerting.py     & уровни severity ladder; пропуск под
порогом; дедупликация в окне; размер истории \\
\hline
\end{tabularx}
\end{table}

\subsection{Соответствие реализации функциональным требованиям}

В таблице~\ref{tab:traceability} приведено сводное соответствие
функциональных требований (FR-1--FR-10) реализованным артефактам и
тестовому покрытию.

\begin{table}[H]
\caption{Соответствие реализации функциональным требованиям}
\label{tab:traceability}
\centering
\footnotesize
\setlength{\tabcolsep}{4pt}
\begin{tabularx}{\linewidth}{|L{2.5cm}|Z|L{2.8cm}|}
\hline
Требование & Артефакт реализации & Тест \\
\hline
FR-1: схема UNSW-NB15      & Pydantic-схема schema.py; валидация в load.py & test\_configs \\
\hline
FR-2: препроцессинг + окна & Preprocessor, build\_windows & test\_preprocess, test\_windowing \\
\hline
FR-3: 15 моделей           & 10 DL + 5 classical в едином registry & test\_models\_smoke \\
\hline
FR-4: порог + калибровка   & TemperatureScaler; find\_threshold\_for\_target\_fpr & test\_metrics \\
\hline
FR-5: формирование алертов & AlertFormer с severity ladder и temporal dedup & test\_alerting \\
\hline
FR-6: REST-API             & FastAPI-сервис inference/service.py & ручное smoke \\
\hline
FR-7: визуальная панель    & Streamlit-дэшборд ui/streamlit\_app.py & ручное smoke \\
\hline
FR-8: FSM-агент            & FSMAgent, DriftInjector, AttackerRuntime & test\_attacker \\
\hline
FR-9: drift-monitoring     & функция drift\_report (PSI, KL, MMD) & test\_drift \\
\hline
FR-10: seed + журнал       & set\_seed; JSON-логи в experiments/ & test\_utils \\
\hline
\end{tabularx}
\end{table}

\noindent Выводы по главе 3

\begin{enumerate}
    \item Реализован программный прототип \texttt{diploma\_nids}
    объемом порядка 4000 строк Python с восемью подпакетами и
    21 конфигурационным YAML-файлом.
    \item Реализованы 15 моделей-кандидатов (10 нейросетевых и 5
    классических) в едином registry с унифицированным интерфейсом.
    \item Реализован конечно-автоматный агент в двойном режиме
    (эмпирический и параметрический), поддерживающий offline и
    online сценарии и обеспечивающий воспроизводимую ground-truth
    по 11 состояниям.
    \item Реализованы подсистемы online-инференса (FastAPI),
    алертирования с severity ladder и дедупликацией, и визуальной
    панели (Streamlit), пригодные для эксплуатации в составе SOC.
    \item Покрытие тестами включает 11 модулей и 86 проходящих
    тестов; критические инварианты pipeline проверены формально.
\end{enumerate}
